% Original analytic teaching diagram, 2026-11-12.
% Chain example of Section 7.33.6 at K = 5, theta = 1, rho = 0.8, drawn on a line at 40 mm per unit:
% target gradient 2 (blue), mean of the target reparameterization formula under the K-step law 2 m_5 = 1.345 (green ring),
% derivative of the K-step program 2(1 - 0.8^5)^2 = 0.904 (yellow).  Small green dots: ten seeded realizations of the
% n = 20 sample average of 2X^(5) (seed 733), illustrating the finite-sample spread around 1.345.
\definecolor{VB}{HTML}{7188A3}\definecolor{VBF}{HTML}{E5EAF0}
\definecolor{VG}{HTML}{829C8B}\definecolor{VGF}{HTML}{E8EEE9}
\definecolor{VY}{HTML}{B4A16F}\definecolor{VYF}{HTML}{F1ECDD}
\definecolor{VR}{HTML}{A7787D}\definecolor{VRF}{HTML}{F0E5E6}
\definecolor{VP}{HTML}{9588A5}\definecolor{VPF}{HTML}{ECE8F0}
\definecolor{VI}{HTML}{3B4252}\definecolor{VM}{HTML}{6B7280}\definecolor{VL}{HTML}{4C4D4F}
\begin{tikzpicture}[x=40mm,y=1mm,>=Stealth,font=\figurefont\fontsize{8}{10}\selectfont,
lab/.style={text=VM,font=\figurefont\fontsize{7.5}{9}\selectfont}]
\draw[VL,line width=.5pt,->] (0.6,0) -- (2.35,0);
% finite-sample realizations
\foreach \v in {1.273,0.812,1.431,1.090,0.629,1.195,1.795,0.925,1.644,1.388} {\fill[VG] (\v,3.2) circle (.55mm);}
\fill[VB] (2,0) circle (1.2mm);
\draw[VG,line width=.9pt,fill=white] (1.345,0) circle (1.2mm);
\fill[VY] (0.904,0) circle (1.2mm);
\node[text=VB,anchor=north] at (2,-1.8) {$\nabla J(\theta)$};
\node[text=VY,anchor=north] at (0.904,-1.8) {$\nabla J_{K}(\theta)$};
\draw[VR,line width=.8pt,<->] (1.345,-8.6) -- (2,-8.6);
\node[lab,text=VR,anchor=north] at (1.67,-9.4) {采样偏差};
\draw[VG,line width=.8pt,<->] (0.629,6.4) -- (1.795,6.4);
\node[lab,text=VG,anchor=south] at (1.345,7) {有限样本波动};
\node[text=VG,anchor=north] at (1.345,-1.8) {$\mathbb{E}[2X^{(K)}]$};
\end{tikzpicture}
